\documentclass[10pt,a4paper]{amsart}
\usepackage[margin=19mm]{geometry}
\usepackage{amsmath,amssymb,mathtools}
\usepackage[scr=rsfs,cal=euler]{mathalfa}
\usepackage{xfrac}
\usepackage[unicode,hidelinks,bookmarks=false]{hyperref}
\newcommand{\ie}{i.e.\ }
\newcommand{\cf}{cf.\ }
\newcommand{\resp}{respectively\ }
\newcommand{\Subsection}{\subsection}
\providecommand{\nobreakdash}{\nobreak\hbox{-}}
\setlength{\emergencystretch}{2em}
\multlinegap0pt
\usepackage{bm}
\usepackage{bbm}
\usepackage{dsfont}
\usepackage{xspace}
\usepackage{cancel}
\usepackage{mathtools}
\usepackage{amsthm}
\makeatletter
\@ifundefined{thm}{\newtheorem{thm}{Theorem}}{}
\@ifundefined{lem}{\newtheorem{lem}[thm]{Lemma}}{}
\makeatother
\title[Weak derivatives and metric differentiability almost everywh]{Weak derivatives and metric differentiability almost everywhere}
\author{Nikita Evseev}
\date{Source: arXiv:2511.02520v1 (CC BY 4.0)}
\begin{document}
\maketitle

\noindent\emph{Source and licence.} Weak derivatives and metric differentiability almost everywhere. Version arXiv:2511.02520v1 (CC BY 4.0), distributed under the Creative Commons Attribution 4.0 International licence, as recorded in the arXiv licence metadata for this exact version and shown on the abstract page https://arxiv.org/abs/2511.02520v1. Full rights notice: licenses/arxiv-2511.02520-CC-BY-4.0.txt.\par\smallskip

\noindent\textit{Editorial scope.} Counted unit: the Thm whose source label is \texttt{theorem:metricSob} in the article (source0.tex lines 807--815, source label \texttt{theorem:metricSob}), delivered together with the article's own complete original proof of it (source0.tex lines 816--888, one \texttt{proof} environment of 73 lines). No mathematics is written, repaired or re-derived: statement and proof are the article's own text line for line. Required context retained uncounted: eq:comisometric (lines 313--315); lemma:lipapprox (lines 621--627); lemma:LipDiff (lines 760--763). Every definition, display and label of those lines is retained unchanged. Omitted from this package and not redistributed: 1--312, 316--620, 628--759, 764--806, 889--996. Nothing inside a retained excerpt is omitted or altered: every retained body is delivered exactly as the article's lines are written, so no omission receipt of any kind accompanies this package. Citations: the retained excerpts carry no citation command at all, so no bibliography entry of the article is transcribed and no reference is redistributed. Reference adaptation: none was needed and none was made; every reference inside the delivered unit resolves inside it, so no unresolved reference or citation is produced. Printed slips kept literal: no printed word, symbol, hypothesis or number of the counted statement or of its proof was altered. Layout and class: the article's own class is \texttt{amsart} with packages including inputenc, stackengine, amssymb, amsthm, amsmath, enumerate, mathtools, comment, todonotes, xcolor, hyperref, dsfont; the wrapper is the lane's pinned \texttt{amsart} editorial wrapper at 10pt on A4, which restates the theorem-like environments the retained text needs and adds the article's own macro definitions that the retained text needs (none), with extra wrapper packages (bm, bbm, dsfont, xspace, cancel, mathtools). The delivered numbering is the unit's own. Assets: no retained span contains a figure environment, a graphics-inclusion command, a PDF-page inclusion, a table or a TikZ picture; no image file is copied into this package, the package declares no asset dependency and no image file is redistributed with it. Licence: version arXiv:2511.02520v1 of this article is distributed under the Creative Commons Attribution 4.0 International licence, as recorded in the arXiv licence metadata for this exact version and shown on the abstract page https://arxiv.org/abs/2511.02520v1. A full rights notice is delivered at licenses/arxiv-2511.02520-CC-BY-4.0.txt. Characters: every retained excerpt is pure ASCII, so the delivered engine, XeLaTeX with its default Latin Modern text font, reports no missing character.
\begin{equation}\label{eq:comisometric}
  \langle \varphi, \partial^\circ_\nu (\psi\circ f)(x) \rangle =   \langle \varphi\circ\psi, \partial^\circ_\nu f(x) \rangle.
\end{equation}

\begin{lem}\label{lemma:lipapprox}
Let $f:\Omega\to V$ be a map belonging to the Sobolev space $W^{1,p}(\Omega; V)$, $p\in [1,\infty)$.
Then for every $\varepsilon>0$ there exists a Lipschitz continuous map $F:\Omega\to V$ such that
\[
|\{x\in\Omega \colon f(x) \ne F(x)\} | < \varepsilon \quad\text{ and }\quad \|f - F \|_{W^{1,p}(\Omega; V)}<\varepsilon.
\]
\end{lem}

\begin{lem}\label{lemma:LipDiff}
Let $f:\Omega\to X$ be a Lipschitz continuous map.
Then, whenever it exists, $\textnormal{md}(f, x)$ is also a metric differential in the topology of $W^{1,p}$, $p\in[1,\infty)$.
\end{lem}

\begin{thm}\label{theorem:metricSob}
Every map in the Sobolev space 
$W^{1,p}(\Omega; X)$, $p\in [1, \infty)$
is almost everywhere metrically differentiable in the topology of $W^{1,p}$.
There exists a seminorm $\rho$ on $\left(\textnormal{Lip}_{z_0}(X)\right)^*$ such that  
  $\mathbb R^n \ni \nu\mapsto \rho(\nabla^of(x)\cdot\nu ) $
is the metric differential of a map $f$
at a point $x\in\Omega$ in the topology of $W^{1,p}$. 
\end{thm}

\begin{proof}
Let $k:X\to V$ be an isometric embedding with $k(z_0)=0$. Then $k\circ f \in W^{1,p}(\Omega; V)$.
Let $\Sigma$ be a null set such that the image $k\circ f(\Omega\setminus\Sigma)$ is separable in $V$ and let $\{\phi_k\}_{k\in\mathbb N}$ be an essential gauge sequence for 
$k\circ f(\Omega\setminus\Sigma)$.
Define seminorm
$
\rho(w) = \sup_k|  \langle \varphi_k, w \rangle|
$
for $w\in \left(\textnormal{Lip}_{0}(V)\right)^*$.

By Lemma \ref{lemma:lipapprox} for each $\varepsilon>0$ there is a measurable set $A_\varepsilon\subset\Omega$
and a Lipschitz continuous map $F_\varepsilon:\Omega\to X$ such that $|\Omega\setminus A_\varepsilon|<\varepsilon$
and $F_\varepsilon(x) = k\circ f(x)$ for $x\in A_\varepsilon$. 
Let $x\in A_\varepsilon$ be a density point, a point of metric differentiability of $F_\varepsilon$,
and a point where partial derivatives $\frac{\partial \phi_k\circ k\circ  f}{\partial x_j}(x)$ and $\frac{\partial \phi_k\circ F_\varepsilon}{\partial x_j}(x)$
exist for $j=1,\dots n$ and $k\in \mathbb N$.
Then $\langle \phi_k, \nabla^{\circ} (k\circ f)(x) \rangle = \langle \phi_k, \nabla^{\circ}F_\varepsilon(x) \rangle $,
and in particular, $\rho(\nabla^{\circ}k\circ f(x)\cdot\nu ) = \rho(\nabla^{\circ}F_\varepsilon (x)\cdot\nu )$.
For $h>0$ denote $D_h = \{\nu\in B \colon x + h\nu \in A_\varepsilon\}$.
Because $x$ is a density point, we have 
\begin{equation}\label{eq:BD-0}
\lim_{h\to 0} |B\setminus D_h| = 0.
\end{equation}
Consider functions
\[
\eta_h(\nu) = \frac{\|k\circ f(x+h\nu) - k\circ f(x)\|_V}{h} - \rho(\nabla^{\circ}(k\circ f)(x)\cdot\nu )
\]
and
\[
\psi_h(\nu) = \frac{\| F_\varepsilon(x+h\nu) -  F_\varepsilon(x)\|_V}{h} - \rho(\nabla^{\circ}F_\varepsilon (x)\cdot\nu ).
\]
Note that $\eta_h(\nu) = \psi_h(\nu)$ for $\nu\in\ D_h$, and
due to Lemma \ref{lemma:LipDiff} 
\begin{equation}\label{eq:psi-0}
\lim_{h\to 0} \| \psi_h\|_{W^{1,p}(B)} = 0.
\end{equation} 
Now for $\varepsilon_1>0$ find Lipschitz function $\xi_h:B\to \mathbb R$ such that 
$\eta_h(\nu) = \xi_h(\nu)$ on measurable set $E_h\subset B$, $|B\setminus E_h| < \varepsilon_1$, and
$\|\eta_h - \xi_h \|_{W^{1,p}(B)}\leq\varepsilon_1$ (Lemma \ref{lemma:lipapprox}).
Then $\psi_h(\nu) = \xi_h(\nu)$ for all $\nu\in D_h\cap E_h$ and $\nabla\psi_h(\nu) = \nabla\xi_h(\nu)$ for almost every $\nu\in D_h\cap E_h$.
Show that the Lipschitz constants $\operatorname{Lip}(\psi_h)$ and $\operatorname{Lip}(\xi_h)$ do not depend on $h$.
Indeed
\[
\operatorname{Lip}(\psi_h) \leq \operatorname{Lip}(F_\varepsilon) + \bigg(\sum_{j=1}^n \|\partial^o_jf(x) \|^2_{\left(\textnormal{Lip}_{0}(X)\right)^*} \bigg)^{\frac{1}{2}}.
\]
While $\operatorname{Lip}(\xi_h)$ depends only on $|\nabla \eta_h (\nu)|$, for which we have the following estimate
\[
|\nabla \eta_h(\nu)| \leq g(x + h\nu) + \bigg(\sum_{j=1}^n \|\partial^o_jf(x) \|^2_{\left(\textnormal{Lip}_{0}(X)\right)^*} \bigg)^{\frac{1}{2}}.
\]
For fixed $\varepsilon_1$ it implies $\operatorname{Lip}(\xi_{h_1}) \leq \operatorname{Lip}(\xi_{h_2})$, when $h_1<h_2$.
So, we can assume that $\operatorname{Lip}(\xi_h) \leq \operatorname{Lip}(\xi_{h_0})$ for all $h\leq h_0$.
To estimate the norm $\|\xi_h - \psi_h \|_{W^{1,p}(B)}$ we note that $B\subset (B\setminus E_h)\cup (B\setminus D_h)\cup (E_h\cap D_h)$
and derive
\begin{align*}
\|\xi_h - \psi_h \|_{W^{1,p}(B)} \leq &C\cdot\varepsilon_1^{\frac{1}{p}} + C|B\setminus D_h|^{\frac{1}{p}} + 0\\
+ &C'\cdot\varepsilon_1^{\frac{1}{p}} + C'|B\setminus D_h|^{\frac{1}{p}} + 0,
\end{align*}
where $C = \sup_{B}|\xi_h(\nu) - \psi_h(\nu)| $ and $C' = 2\max\{\operatorname{Lip}(\psi_h), \operatorname{Lip}(\xi_h) \}$,
both constants do not depend on $h$.
Finally, from \eqref{eq:BD-0}, \eqref{eq:psi-0} and 
\[
\|\eta_h \|_{W^{1,p}(B)} \leq \|\eta_h -\xi_h \|_{W^{1,p}(B)} + \|\xi_h - \psi_h \|_{W^{1,p}(B)} + \|\psi_h \|_{W^{1,p}(B)}
\]
we have
\[
\limsup_{h\to 0 } \|\eta_h \|_{W^{1,p}(B)} \leq \varepsilon_1 + (C+ C')\varepsilon_1^{\frac{1}{p}}.
\]
Due to the arbitrariness of $\varepsilon_1$ and $\varepsilon$, we conclude that $k\circ f$ is almost everywhere metrically differentiable in the topology of $W^{1,p}$
with metric differential $\nu\mapsto \sup_k|  \langle \varphi_k, \nabla^{\circ} (k\circ f)(x)\cdot\nu \rangle|$.
It remains to note that $d(f(x+h\nu), f(x)) = \|k\circ f(x+h\nu) - k\circ f(x)\|_V$. Therefore mapping $f$ is also almost everywhere metrically differentiable in the topology of $W^{1,p}$
and with property \eqref{eq:comisometric} its corresponding metric differential could be writen as 
$\nu\mapsto \sup_k|  \langle \phi_k\circ k, \nabla^{\circ} f(x)\cdot\nu \rangle|$.
\end{proof}

\end{document}
